\section{Conclusion}\label{sec:conclusion}

In this paper, I measure club-specific refereeing gaps across the top-5 European leagues and the Champions League over a quarter century and place FC Barcelona in that distribution. Barcelona's raw advantage in cards and fouls is large, but most of it reflects the way Barcelona plays. Conditional on fouls called, referees do not book Barcelona's players less or its opponents more. On penalties, Barcelona's treatment is unremarkable. The decline in Barcelona's card advantage after 2017/18 is real, but its timing tracks the club's sporting decline at least as closely as the end of the reported payments. The clearest favoritism in the data, extra stoppage time when trailing, is shared equally by Real Madrid.

These results bound, but do not settle, the question of whether the payments affected outcomes. Aggregate decision counts cannot detect favoritism concentrated in a few high-stakes decisions, in particular referees' assignments, or in decisions not recorded in box scores. Three extensions would sharpen the analysis: referee-level identities for La Liga after 2011/12, minute-by-minute card timing by score state, and VAR interventions by team since 2018/19.
